\documentclass[12pt, a4paper]{report}

% Packages :

\usepackage[french]{babel}
\usepackage[utf8]{inputenc}
\usepackage[T1]{fontenc}
\usepackage[pdftex, pdfauthor={Bacomathiques}]{hyperref}
\usepackage{sectsty}
\usepackage[explicit]{titlesec}
\usepackage{xcolor}
\usepackage{amsmath}
\usepackage{amssymb}
\usepackage{amsthm}
\usepackage{fourier}
\usepackage{titlesec}
\usepackage{fancyhdr}
\usepackage{catchfilebetweentags}
\usepackage[french, capitalise, noabbrev]{cleveref}
\usepackage[fit, breakall]{truncate}
\usepackage[margin=3cm]{geometry}
\usepackage{tocloft}
\usepackage{tikz}
\usepackage{tocloft}
\usepackage{microtype}
\usepackage{listings}
\usepackage{tabularx}
\usepackage{calc}
\usepackage[export]{adjustbox}
\usepackage[most]{tcolorbox}
\usepackage{standalone}
\usepackage{xlop}
\usepackage{etoolbox}
\usepackage{environ}

\usetikzlibrary{arrows.meta}
\usetikzlibrary{trees}

% Paramètres :

\author{Bacomathiques}
\definecolor{graphe}{HTML}{93c9ff}
\setcounter{MaxMatrixCols}{12}
\setlength{\parindent}{0pt}
\setlength{\fboxsep}{0pt}
%\pdfsuppresswarningpagegroup=1

% Code :

\lstdefinestyle{style}{
	backgroundcolor=\color{white},
	commentstyle=\em\color[HTML]{999988},
	keywordstyle=\bfseries,
	identifierstyle=\normalfont,
	stringstyle=\color[rgb]{0.87, 0.07, 0.27},
	basicstyle=\ttfamily\color{black},
	breakatwhitespace=false,
	breaklines=true,
	captionpos=b,
	keepspaces=true,
	numbers=left,
	numbersep=5pt,
	showspaces=false,
	showstringspaces=false,
	showtabs=false,
	tabsize=2,
	numbers=none
}

\lstset{style=style}
\lstset{
	literate=
	{á}{{\'a}}1
	{à}{{\`a}}1
	{ã}{{\~a}}1
	{é}{{\'e}}1
	{ê}{{\^e}}1
	{í}{{\'i}}1
	{ó}{{\'o}}1
	{õ}{{\~o}}1
	{ú}{{\'u}}1
	{ü}{{\"u}}1
	{ç}{{\c{c}}}1
}

\lstset{
	framextopmargin=10pt,
	framexrightmargin=10pt,
	framexbottommargin=10pt,
	framexleftmargin=10pt,
	xleftmargin=10pt,
	xrightmargin=10pt,
}

% Couleurs :

\definecolor{title}{HTML}{912c21}
\definecolor{section}{HTML}{1c567d}
\definecolor{subsection}{HTML}{2980b9}

\definecolor{rule}{HTML}{c4c4c4}

\definecolor{formula}{HTML}{ebf3fb}
\definecolor{formula-left}{HTML}{3583d6}

\definecolor{tip}{HTML}{dcf3d8}
\definecolor{tip-left}{HTML}{26a65b}

\definecolor{demonstration}{HTML}{fff8de}
\definecolor{demonstration-left}{HTML}{f1c40f}

\definecolor{exercise}{HTML}{e0f2f1}
\definecolor{exercise-left}{HTML}{009688}

\definecolor{correction}{HTML}{e0f7fa}
\definecolor{correction-left}{HTML}{00bcd4}

\definecolor{toc}{HTML}{fceae9}
\definecolor{toc-left}{HTML}{e74c3c}
\definecolor{toc-dark}{HTML}{87281f}

% Titres :

\renewcommand{\thesection}{\Roman{section} - }
\renewcommand{\thesubsection}{\arabic{subsection}. }

\newcommand{\sectionstyle}{\normalfont\LARGE\bfseries\color{section}}
\titleformat{\section}{\sectionstyle}{\thesection #1}{0pt}{}
\titleformat{name=\section, numberless}{\sectionstyle}{#1}{0pt}{}

\newcommand{\subsectionstyle}{\normalfont\Large\bfseries\color{subsection}}
\titleformat{\subsection}{\subsectionstyle}{\thesubsection #1}{0pt}{}
\titleformat{name=\subsection, numberless}{\subsectionstyle}{#1}{0pt}{}

\titlelabel{\thetitle\ }

% Table des matières :

\addto\captionsfrench{\renewcommand\contentsname{}}
\renewcommand{\cftsecpagefont}{\color{toc-dark}}
\renewcommand{\cftsubsecpagefont}{\color{toc-dark}}
\renewcommand{\cftsecleader}{\cftdotfill{\cftdotsep}}
\renewcommand{\cftsecfont}{\bfseries}
\renewcommand{\cftsecpagefont}{\bfseries\color{toc-dark}}
\setlength{\cftbeforetoctitleskip}{0pt}
\setlength{\cftaftertoctitleskip}{0pt}
\setlength{\cftsecindent}{0pt}
\setlength{\cftsubsecindent}{20pt}
\setlength{\cftsubsecnumwidth}{20pt}

% Commandes :

\newcommand{\newpar}{\\[\medskipamount]}
\newcommand{\lesson}[3]{%
	\newcommand{\level}{#1}%
	\newcommand{\id}{#2}%
	\hypersetup{pdftitle={#3}}
	\begin{center}%
		\includegraphics[width=150px]{\imagespath/bacomathiques}%
		
		\vspace{30pt}%
		{\Huge\color{title} #3}%
		
		\vspace{10pt}%
		{Bacomathiques --- \href{https://bacomathiqu.es/cours/#1/#2}{\color{section} https://bacomathiqu.es}}%
		
		\vspace{20pt}%
	\end{center}%
	\begin{toc}
		\tableofcontents%
	\end{toc}
	\thispagestyle{empty}%
	\newpage%
	\setcounter{page}{1}%
}
\newcommand{\imagespath}{../../images}
\newcommand{\lessonimagespath}{\imagespath/lessons/\level/\id/}
\newcommand{\includelatexpicture}[2][\textwidth - 100pt]{%
	\begin{center}%
		\resizebox{#1}{!}{%
			\input{\lessonimagespath#2}%
		}%
	\end{center}%
	\medskip%
}
\newcommand{\includeimage}[1]{%
	\begin{center}%
		\includegraphics{\lessonimagespath#1}%
	\end{center}%
	\medskip%
}
\newcommand{\includerepresentation}[1]{%
	\begin{center}%
		\setlength{\fboxrule}{0.5pt}%
		\href{https://www.geogebra.org/m/#1}{\includegraphics[width=\textwidth-1pt,fbox]{\lessonimagespath#1}}%
	\end{center}%
}
\newcommand{\floor}[1]{\lfloor #1 \rfloor}

\makeatletter
\newcommand\inputcontent{\@ifstar{\inputcontent@star}{\inputcontent@nostar}}
\newcommand{\inputcontent@star}[1]{%
	\ExecuteMetaData[#1]{content}%
}
\newcommand{\inputcontent@nostar}[1]{%
	\newpage%
	\inputcontent@star{#1}%
}
\makeatother

\let\oldsection\section
\renewcommand\section{\clearpage\oldsection}
\newcommand{\contentwidth}[1][medium]{}

% En-têtes :

\pagestyle{fancy}

\renewcommand{\sectionmark}[1]{\markboth{\thesection \ #1}{}}

\fancyhead[R]{\truncate{0.23\textwidth}{\color{title}\thepage}}
\fancyhead[L]{\truncate{0.73\textwidth}{\color{title}\leftmark}}
\fancyfoot[C]{\scriptsize \href{https://bacomathiqu.es/cours/\level/\id}{\texttt{bacomathiqu.es}}}

\makeatletter
\patchcmd{\f@nch@head}{\rlap}{\color{rule}\rlap}{}{}
\patchcmd{\headrule}{\hrule}{\color{rule}\hrule}{}{}
\makeatother

% Environnements :

\newenvironment{nosummary}{}{}
\newcommand{\tcolorboxtitle}[2]{\setlength{\fboxsep}{2.5pt}\hspace{-10pt}\colorbox{#1-left}{\hspace{8pt}\MakeUppercase{#2} \hspace{2pt} \includegraphics[height=0.8em]{\imagespath/bubbles/#1}\hspace{5pt}}}
\newcommand{\tcolorboxsubtitle}[2]{\ifstrempty{#2}{}{\textcolor{#1-left}{\large#2}\\[\medskipamount]}}
\tcbset{
	frame hidden,
	boxrule=0pt,
	boxsep=0pt,
	enlarge bottom by=8.5pt,
	enhanced jigsaw,
	boxed title style={sharp corners,boxrule=0pt,coltitle={white},titlerule=0pt},
	fonttitle=\fontsize{6pt}{6pt}\bfseries\boldmath,
	top=10pt,
	right=10pt,
	bottom=10pt,
	left=10pt,
	arc=0pt,
	outer arc=0pt,
}
\newtcolorbox{toc}[1][]{
	colback=toc,
	borderline west={3pt}{0pt}{toc-left},
	title=\tcolorboxtitle{toc}{Table des matières},
	colbacktitle=toc,
	before upper={\tcolorboxsubtitle{toc}{#1}}
}
\newtcolorbox{formula}[1][]{
	colback=formula,
	borderline west={3pt}{0pt}{formula-left},
	title=\tcolorboxtitle{formula}{À retenir},
	colbacktitle=formula,
	before upper={\tcolorboxsubtitle{formula}{#1}}
}
\newtcolorbox{tip}[1][]{
	colback=tip,
	borderline west={3pt}{0pt}{tip-left},
	title=\tcolorboxtitle{tip}{À lire},
	colbacktitle=tip,
	before upper={\tcolorboxsubtitle{tip}{#1}}
}
\newtcolorbox{demonstration}[1][]{
	colback=demonstration,
	borderline west={3pt}{0pt}{demonstration-left},
	title=\tcolorboxtitle{demonstration}{Démonstration},
	colbacktitle=demonstration,
	before upper={\tcolorboxsubtitle{demonstration}{#1}}
}

\NewEnviron{whitetabularx}[1]{%
	\renewcommand{\arraystretch}{2.5}
	\colorbox{white}{%
		\begin{tabularx}{\textwidth}{#1}%
			\BODY%
		\end{tabularx}%
	}%
}

% Longueurs :

\newlength{\espacetitreliste}
\setlength{\espacetitreliste}{-16pt}
\newcommand{\entretitreetliste}{\vspace{\espacetitreliste}}

\begin{document}
	%<*content>
	\lesson{terminale}{5}{fonction-logarithme}{La fonction logarithme népérien}
	
	\header{caption}{Le logarithme népérien est utilisé dans les calculs logarithmiques et dans
		l'étude de certaines croissances.}
	
	\header{excerpt}{Une fonction logarithme est une fonction possédant la propriété de pouvoir
		transformer des produits en sommes. L''utilisation de telles fonctions permet de
		faciliter les calculs comprenant de nombreuses multiplications, divisions et élévations
		à des puissances rationnelles. Ici, nous étudierons la fonction logarithme népérien
		: relations algébriques, signe, dérivée, variations, représentation et plus encore
		!}
	
	\header{difficulty}{4}

	\section{Propriétés du logarithme népérien}

	\subsection{Définition}

	\begin{formula}[Définition]
		Le \textbf{logarithme népérien} notée $\ln$ est la fonction définie sur $]0;+\infty[$ telle que pour tout $x > 0$ et $y$ réels :
		\[ \ln(x) = y \iff x = e^y \]
	\end{formula}

	Ainsi, a tout réel \textbf{strictement positif} $x$, la fonction logarithme népérien y associe \textbf{son unique antécédent} $y$ par rapport à \href{https://bacomathiqu.es/cours/premiere/fonction-exponentielle/}{la fonction exponentielle}. De même pour la fonction exponentielle.

	On dit que ces fonctions sont des \textbf{fonctions réciproques} (à la manière de $\sin$ et $\arcsin$ ou $\cos$ et $\arccos$).

	\begin{tip}[Exemple]
		Cette relation peut sembler compliquer à assimiler mais il n'en est rien ! Prenons $x = 0$, on a :
		\newpar
		$e^0 = 1$ (tout réel mis à la puissance zéro vaut un), la relation précédente nous donne $\ln(1) = 0$.
		\newpar
		Si on prend maintenant $x = 1$, on a :
		\newpar
		$e^1 = e$, on a donc $\ln(e) = 1$.
	\end{tip}

	Les relations suivantes sont par conséquent disponibles :

	\begin{formula}[Relations entre fonctions réciproques]
		Pour tout réel $x$ \textbf{strictement positif}, on a $e^{\ln(x)} = x$.
		\newpar
		Et pour tout réel $x$, on a $\ln(e^x) = x$.
	\end{formula}

	\subsection{Relations algébriques}

	Le logarithme népérien a plusieurs propriétés intéressantes qu'il faut connaître.

	\begin{formula}[Formules]
		Pour tous réels $x$ et $y$ \textbf{strictement positifs} :
		\begin{itemize}
			\item $\ln(x \times y) = \ln(x) + \ln(y)$
			\item $\ln(x^n) = n \times \ln(x)$ pour $n \in \mathbb{Z}$
			\item $\ln\left(\frac{x}{y}\right) = \ln(x) - \ln(y)$
			\item $\ln\left(\frac{1}{y}\right) = -\ln(y)$
			\item $\ln(\sqrt[p]{x}) = \frac{1}{p} \times \ln(x)$ pour $p \in \mathbb{N}^*$
		\end{itemize}
	\end{formula}

	Certaines de ces propriétés peuvent se déduire les unes des autres.

	\subsection{Représentation graphique}
	\label{representation-graphique}

	Voici une représentation graphique de la fonction logarithme népérien :

	\includerepresentation{wsntfeab}

	On voit sur ce graphique plusieurs propriétés données précédemment : $\ln(1) = 0$ et $\ln(e) = 1$ par exemple.
	On trace maintenant le graphe de la fonction logarithme népérien, avec celui de la fonction exponentielle. On trace également la droite d'équation $y = x$ :

	\includerepresentation{aymmb94w}

	On remarque plusieurs choses : le graphe de la fonction logarithme népérien est le symétrique de celui de la fonction exponentielle par rapport à la
	droite $y = x$ et on voit que la fonction logarithme népérien croît moins vite que la fonction puissance qui elle-même croît moins vite que la fonction exponentielle. Cette propriété est importante : c'est la \textbf{croissance comparée}.

	\section{Étude de la fonction}

	\subsection{Limites}

	\begin{formula}[Limites]
		Les limites de la fonction logarithme népérien aux bornes de son ensemble de définition sont :
		\begin{itemize}
			\item $\lim\limits_{x \rightarrow 0^+} \ln(x) = -\infty$
			\item $\lim\limits_{x \rightarrow +\infty} \ln(x) = +\infty$
		\end{itemize}
	\end{formula}

	Il faut aussi savoir que la fonction puissance ``l'emporte'' sur le logarithme népérien (voir la partie \hyperref[representation-graphique]{``Représentation graphique''}).

	\begin{formula}[Croissances comparées]
		Pour tout $n \in \mathbb{N}$ :
		\begin{itemize}
			\item $\lim\limits_{x \rightarrow +\infty} \frac{\ln(x)}{x^n} = 0$.
			\item $\lim\limits_{x \rightarrow 0^+} x^n \ln(x) = 0$.
		\end{itemize}
	\end{formula}

	\begin{demonstration}[Croissances comparées]
		Nous allons démontrer le second point en utilisant le premier (qui n'est pas éligible à une démonstration au lycée) dans le cas $n = 1$. Pour tout $x > 0$, posons $y = \frac{1}{x}$.
		\newpar
		On a donc pour tout, $x \ln(x) = \frac{1}{y} \ln\left(\frac{1}{y}\right) = -\frac{\ln(y)}{y}$.
		\newpar
		Or, quand $x$ tend vers $0^+$, $y$ tend vers $+\infty$. Par le premier point :
		\[ \lim\limits_{\substack{y \rightarrow +\infty}} \frac{\ln(y)}{y} = 0 \iff \lim\limits_{\substack{y \rightarrow +\infty}} -\frac{\ln(y)}{y} = 0 \]
		Et en remplaçant $y$ par $\frac{1}{x}$ dans le résultat ci-dessus, on a bien ce que l'on cherchait.
	\end{demonstration}

	Pour finir, on donne une limite qu'il peut être utile de savoir redémontrer.

	\begin{formula}
		\[ \lim\limits_{\substack{x \rightarrow 0}} \frac{\ln(1 + x)}{x} = 1 \]
	\end{formula}

	\begin{demonstration}
		La fonction logarithme népérien est dérivable en $1$ (voir sous-section suivante), on peut donc écrire :
		\[ \lim\limits_{x \rightarrow 0} \frac{\ln(1+x)}{x} = \lim\limits_{x \rightarrow 0} \frac{\ln(x+1) - \ln(1)}{x} = \ln'(1) = \frac{1}{1} = 1 \]
	\end{demonstration}

	\subsection{Dérivée}

	\begin{formula}[Dérivée d'une composée]
		Soit une fonction $u$ dérivable et \textbf{strictement positive} sur un intervalle $I$, on a pour tout $x$ appartenant à cet intervalle :
		\[ \ln'(u(x)) = \frac{u'(x)}{u(x)} \]
	\end{formula}

	\begin{formula}[Dérivée]
		Ainsi, si pour tout $x \in I$ on a $u(x) = x$, on trouve :
		\[ \ln'(x) = \frac{1}{x} \]
	\end{formula}

	\subsection{Variations}

	Avec la dérivée donnée précédemment ainsi que les limites données, il est désormais possible d'obtenir les variations de la fonction logarithme népérien.

	\begin{formula}[Signe et variations]
		\contentwidth[big]
		\includelatexpicture{variations}

		On remarque qu'avec le tableau de variation, il est possible d'obtenir le signe de la fonction (avec le théorème des valeurs
		intermédiaires).
		\newpar
		Ainsi, sur $]0;1[$, $\ln$ est \textbf{strictement négative} et sur $]1;+\infty[$,
		$\ln(x)$ est \textbf{strictement positive} et, comme vu précédemment, $\ln(1) = 0$.
		\newpar
		On observe également les variations de la fonction : strictement croissante sur son ensemble de définition.
	\end{formula}
	%</content>
\end{document}
